The preceding rank and lifting statements concern the symbol
with the reduced weights. In Proposition~\ref{prop:residue-plane-neighbor-test}
the source must have the full weight. The replacement below
retains the rank of the model family and supplies the integral
contractions used on the target side.

\paragraph{Restoring the bounded weights.}
\label{par:plane-middle-weights}
\hypertarget{source-plane-middle-weights}{}
At a top let \(\lambda=(\lambda_1,\ldots,\lambda_k)\) be its
true dominant weight and let \(b=\lambda_1-\lambda_2\). Remove
\((b,0,\ldots,0)\) to form the remaining dominant weight
\[
 \mu=(\lambda_2,\lambda_2,\lambda_3,\ldots,\lambda_k).
\]

Away from the tops take the removed weight to be zero and
\(\mu=\lambda\). At odd indices use the corresponding reversed
negative weights in the polynomial model. Write \(d_l(\gamma)\)
for the dimension of highest weight \(\gamma\) in dimension \(l\).
The local factor of the count for restoring weights is
\(d_{k-1}(\lambda_2,\ldots,\lambda_k)\) at a top and
\(d_k(\lambda)\) elsewhere. We construct integral polynomials
realizing exactly this count.

\begin{auxlemma}[Counting the tableaux at one index]\label{lem:bounded-tableau-count}
With the top weights \(\mu\) defined at the start of
paragraph~\ref{par:plane-middle-weights}, use tableaux of the
determinant-shifted shape. At an even top take tableaux with every entry in the first row
equal to zero. At an odd top take the reversed negative shape and exclude zero
from its alphabet. Writing \(n_s\) for the number admitted at position \(s\),
\[
 n_s=\begin{cases}d_{k-1}(\lambda_2,\ldots,\lambda_k),&s\text{ a top},\\d_k(\lambda),&s\text{ away from the tops}.\end{cases}
\]
Away from the tops every tableau is allowed. Distinct selected tableaux have distinct leading monomials over
every residue field.
\end{auxlemma}
\begin{proof}
Use a matrix of variables \(X=(X_{i,h})\), with rows and columns
numbered from zero. For a dominant weight \(\gamma\) the
polynomial realization after its determinant twist has shape
\(\gamma-\gamma_k(1,\ldots,1)\). Its integral basis consists of
products of initial-column minors indexed by semistandard
tableaux of this shape on the alphabet \(0,\ldots,k-1\).
For each tableau its leading monomial uses the diagonal of each
minor. Order the variables by column first and by increasing
row within a column, giving the earlier variable precedence.

In a minor the largest monomial uses the smallest available row
in the first column and then the smallest remaining row in each
successive column. For rows \(i_0<\cdots<i_{s-1}\), this gives
\[
 \operatorname{in}\det(X_{i_a,b})_{0\le a,b<s}
    =\prod_{b=0}^{s-1}X_{i_b,b}.
\]
The coefficient is one. If \(M_T\) is the diagonal monomial of the product for a tableau \(T\), then
\[
 \deg_{X_{i,b}}M_T=\#\{c:T(b,c)=i\}.
\]
Its weakly
increasing order reconstructs that row from these counts. These are the integral
standard monomial coordinates used for the component lattices.
\cite[Theorem II.2.16]{abw-schur}

Away from the tops take every tableau of \(\mu\) in the
corresponding signed realization. At an even top restrict to
those with first tableau row all zero. Removing that row
leaves the tableau of the smaller shape with alphabet
\(1,\ldots,k-1\), since every remaining box had a box above it.
This restriction counts
\[
 d_{k-1}(\lambda_2,\ldots,\lambda_k),
\]

by subtracting the common determinant shift. At an odd top
restrict to tableaux with no entry zero. The reversed negative
weight of \(\mu\) has two equal last entries. The determinant
shift leaves at least its last shape row empty. Dropping this
row and using the alphabet \(1,\ldots,k-1\) counts the dimension
of the reversed negative of \((\lambda_2,\ldots,\lambda_k)\),
which equals \(d_{k-1}(\lambda_2,\ldots,\lambda_k)\) by duality.
\end{proof}
Let \(n_2\) be the product of these counts and \(H_i\) the
resulting basis products in the independent block variables.

\begin{proposition}[Restoration of the bounded highest weights]\label{prop:restore-plane-model-rank}
Use the removed flat degrees and the \(n_2\) tableau products
\(H_i\) counted in Lemma~\ref{lem:bounded-tableau-count}. The products \(\theta(Q)H_i\) retain
independence modulo the uniformizer for every independent
family of flat polynomials \(Q\) of the prescribed degrees.
With the common factor \(C\) of \eqref{eq:plane-clean-factor}
and precision \(c'\ge\max(2(D_1+1),c_\star+o_S)\), the products
admit the stated clean approximations at the distinguished edge. With the bounded spread \(c_0\) and lower
degree \(a\), let \(f_{\rm true}\) and \(f_{\rm flat}\) be the retained fractions in the true and flat dimensions, respectively. In the indicated size range,
\[
 f_{\rm true}\ge(1-kc_0/a)^N f_{\rm flat}.
\]

\end{proposition}

\begin{proof}
\proofstep{Independence after restoring the removed weight.}
To prove independence after restoring the removed weight, let
\(\theta\) insert the first column at an even index and the
signed \((k-1)\)-row minors of the first \(k-1\) columns at an
odd index. This is the vector and covector polynomial model of
Section~\ref{sec:edge-transfer}. At an even top with removed
degree \(b\), the leading monomial of a degree-\(b\) monomial
in the flat variables inserts exponents \(\alpha_c\) at
\(X_{c,0}\). The exponent matrix of a chosen tableau has no
\(X_{c,0}\) for \(c>0\).

At each \(c>0\) the exponent of \(X_{c,0}\) in the leading product
is therefore precisely \(\alpha_c\). Recover the missing exponent from
the homogeneous degree:
\[
 \alpha_0=b-\sum_{c>0}\alpha_c.
\] At an odd top the leading
monomial of the cofactor omitting row \(c\) uses \(X_{t,t}\)
for \(0\le t<c\) and \(X_{t+1,t}\) for \(c\le t<k-1\).
In the selected tableau the entry in row \(t\) is at least
\(t+1\), so its exponent at \(X_{t,t}\) is zero. The diagonal exponents of the product read
\(d_t=\sum_{c>t}\alpha_c\) for \(0\le t<k-1\).

If \(b\) is the fixed flat degree, invert this triangular
reading by
\[
     \alpha_0=b-d_0,\qquad
     \alpha_c=d_{c-1}-d_c\ (1\le c\le k-2),\qquad
     \alpha_{k-1}=d_{k-2}.
\]
Thus the diagonal exponents and the degree recover the
full flat monomial. After recovering the flat exponents, subtract their contribution
to recover the leading exponent matrix of the tableau. At an
off-top position there is no contribution to subtract. Thus the
pairs of a flat monomial and a selected tableau product have
distinct leading monomials. Each has leading coefficient a unit,
including in the residue field. In a proposed nonzero relation
\[
 \sum_i\theta(Q_i)H_i=0,
\]
choose the largest of those leading
monomials among the nonzero coefficients. Its coefficient cannot
cancel with that of any other pair. This proves the assertion
over the residue field for arbitrary \(Q_i\) of the fixed degrees.

\proofstep{The three normalized edge lattices.}
Use the normalized edge coordinates on all three weights
\(\lambda,\mu,\lambda-\mu\). The normalization consists of
pullback to the component frames, the last-entry determinant
powers in the signed realizations, and the common edge powers
\(w^{A_r}\) for the respective weights. In each frame the common powers on the two factors have product
\[
 w^{A_r(\mu)}w^{A_r(\lambda-\mu)}=w^{A_r(\lambda)},
\]
by additivity of the edge exponent in the weight.

On a flat polynomial this change to edge
coordinates is the letter substitution followed by \(\theta\).
Indeed a vector transforms by the frame matrix \(M\), and the
cofactors by \(\det(M)M^{-t}\), which follows by multiplying
the adjugate identity by \(M\). The last-entry twist on the
covector cancels the determinant. The flat edge exponent is
one copy of the degree for each even block and zero for each
odd block. Write
\[
 D=\operatorname{diag}(w^{-1}I_r,I_{k-r})
\]
for the diagonal
part. After the common determinant cancels, the two letter maps are
\[
 wD=\operatorname{diag}(I_r,wI_{k-r})\quad\text{on vectors},\qquad
 D^{-t}=\operatorname{diag}(wI_r,I_{k-r})\quad\text{on covectors}.
\] The
other frame factors act by the same matrix identities, with
Frobenius at the corresponding index. Thus this normalization
pulls the edge coefficient lattices back to their integral
standard monomial models, including their reductions.

In the distinguished frame the model vectors \(S_b\) have
normalized polynomials \(B_b C\) with one common factor \(C\).
This is \(\operatorname{subst}(\mathrm{SPtrue})\) in the symbol
identity, with residue the product of \(F_*\) and the separate
determinant reserve on the digits. Each determinant uses distinct
rows and has a permutation monomial with unit coefficient, so
\(\overline C\ne0\). To obtain the clean source, expand \(B_b\)
by its segment presentation. Its nonnegative separation margin
allows each homogeneous pair to be expanded in the paired paths
of Section~\ref{par:plane-column-orders}. The common factor
\begin{equation}\label{eq:plane-clean-factor}
                 C=F_*\prod_s T_s^{kQ}
\end{equation}
supplies the precision reserve and the factors for the digits. Write \(S_{b,o}\) for the standard-edge pullback of the
normal degree-\(o\) component of \(C B_b\). Apply the path
and digit constructions to this component.

In the flat source
lattice \(R_e\) on the chosen edge, the resulting vector
\(F_{b,o}\) has source order at most \(o\) and satisfies
\[
                 F_{b,o}-w^o S_{b,o}\in w^{o+c'}R_e .
\]
The insertions commute with
the projection to the full true multidegree because \(C\)
has the reserved degrees. Divide by \(w^o\) in the displayed bound,
then sum over the finitely many \(o\) bounded in
Section~\ref{par:plane-rank}. Thus
\[
  F_b:=\sum_o w^{-o}F_{b,o}\in S_b+w^{c'}R_e .
\]
Here \(S_b=\sum_o S_{b,o}\). The competing-edge conditions
of the construction give \(F_b\in w^{c'}R_{e'}\) on every
other edge out of the same source. Restore the weights by multiplying
with the inverse edge transports of the \(H_i\). They are integral
on the chosen edge, so the same congruence holds in the restored
lattice there.

If a combination with coefficients \(c_{b,i}\in O\) is
zero modulo \(w\) in the normalized edge lattice, then in the
residue polynomial model
\[
 0=\sum_i\theta\left(\overline C\sum_b\overline{c_{b,i}}\overline{B_b}\right)
                    \overline{H_i}.
\]
All polynomials inside \(\theta\) have the full common flat
multidegree. By block independence, for each \(i\) the coefficient polynomial
vanishes. Cancellation takes place in the residue polynomial ring,
a domain, and then independence of the \(\overline{B_b}\) gives
\[
 \overline C\sum_b\overline{c_{b,i}}\overline{B_b}=0
 \ \Longrightarrow\ \sum_b\overline{c_{b,i}}\overline{B_b}=0
 \ \Longrightarrow\ c_{b,i}\in wO\quad\text{for every }b .
\] In particular primitive combinations remain primitive
in the actual source lattice.

\proofstep{Pairing the restored model family.}
For the pairing one can therefore test the restored model family.
An integral target test evaluates the difference between the clean
and model combinations in \(w^{c'}O\). If \(c'\ge 2(D_1+1)\), this
is contained in \(p^{D_1+1}O\), since \(w^2=p\). A model evaluation
of valuation at most \(D_1\) has the same valuation on the clean
combination. On competing source edges restoration costs at most
\(o_S\) powers of \(w\), by the source vertex bound for the remaining
weights. Thus \(c'\ge c_\star+o_S\) gives the source conductor there.

At each top the spread \(\lambda_2-\lambda_j\) for \(j\ge2\)
is at most \(c_0\), with \(b\ge a\) in the plane count.
The Weyl factors involving the first entry give
\[
 d_k(\lambda)\le d_{k-1}(\lambda_2,\ldots,\lambda_k)D_k(b)
                  (1+c_0/a)^{k-1}.
\]
This follows term by term from
\(\lambda_1-\lambda_j+j-1\le (b+j-1)(1+c_0/a)\).
The retained fraction of the true dimension is therefore at
least the flat fraction multiplied by \((1-kc_0/a)^N\), using
\((1+t)^{-(k-1)}\ge1-(k-1)t\) for \(t\ge0\). The product of
the off-top dimensions is included exactly.
\end{proof}
